\begin{textsample}{POS Dim 3 – 2012\_c6 – Score -7.00 – 2012\_c6/enade\_2012\_tec\_gest\_fin\_8526ca4e82.txt}  \label{ex:f3_pos_261}
CAPÍTULO 3 ANÁLISE TÉCNICA DA PROVA Este capítulo tem por objetivo apresentar o desempenho dos estudantes concluintes de Tecnologia em Gestão Financeira no ENADE/2012.
Para isso, foram calculadas as estatísticas básicas da prova em seu todo, bem como as estatísticas dos componentes relacionadas à Formação Geral, ao de Conhecimento Específico da Área e das questões discursivas isoladamente.
Nas tabelas, são apresentados o tamanho da população inscrita e de presentes, e as seguintes estatísticas das notas⁹ : média do desempenho na prova, erro padrão da média, desvio padrão, nota mínima, mediana e nota máxima.
Tais estatísticas contemplam o total de estudantes concluintes da Área de Tecnologia em Gestão Financeira presentes à prova do ENADE 2012, tendo -se em vista as seguintes agregações : ( a ) as Grandes Regiões e o país como um todo ; ( b ) a Categoria Administrativa ; e ( c ) a Organização Acadêmica.
Em relação aos gráficos de distribuição de notas, o intervalo considerado foi de 10 unidades, aberto à esquerda e fechado à direita, com exceção do primeiro intervalo, [ 0 ; 10 ], fechado em ambos os extremos.
Para os gráficos de distribuição das notas das questões discursivas, foram consideradas mais duas categorias : questão em branco e nota zero. 3.1 ESTATÍSTICAS BÁSICAS DA PROVA 3.1.1 Estatísticas Básicas Gerais A Tabela 3.1 apresenta as estatísticas básicas da prova por grande Região.
A população total de inscritos foi de 12.657.
Destes, 9.337 estiveram presentes, sendo 26,2 \% o índice de não comparecimento.
A Região de maior abstenção foi a Sul ( 32,8 \% ) e a de menor abstenção foi a Norte ( 18,2 \% ).
A média das notas da prova como um todo ( nas seções seguintes serão analisados os componentes de Formação Geral e de Conhecimento Específico ) foi 32,8, sendo que os alunos da região Norte obtiveram a média mais baixa ( 30,9 ) e os da região Nordeste obtiveram a média mais alta ( 34,3 ).
As demais médias foram : 33,2 na região Sul, 32,8 na região Centro-Oeste e 32,4 na região Sudeste.
O desvio padrão para o Brasil como um todo foi 11,9, sendo o maior desvio padrão encontrado na região Centro-Oeste ( 13,9 ) e o menor na região Norte ( 10,8 ), indicando uma menor dispersão das notas desta última região. ⁹ Essas estatísticas e outras estão definidas no Capítulo 1. foi 11,9, sendo o maior desvio padrão encontrado na região Centro-Oeste ( 13,9 ) e o menor na região Norte ( 10,8 ), indicando uma menor dispersão das notas desta última região.
A região que obteve a maior nota máxima foi a Sul ( 83,8 ), ao passo que a região que atingiu a menor nota máxima foi a Norte ( 64,6 ).
A mediana do Brasil como um todo foi 31,8, sendo a maior mediana obtida na região Nordeste ( 32,8 ) e a menor obtida na Norte ( 27,4 ).
A nota mínima não foi zero na maioria das regiões, excetuam -se as regiões Norte ( 10,8 ) e Nordeste ( 3,6 ).
O comportamento das notas dos estudantes de todo o Brasil pode ser observado no Gráfico 3.1 que apresenta um histograma com a distribuição das mesmas.
Essa é uma distribuição unimodal com moda no intervalo ( 30 ; 40 ].
O coeficiente de assimetria da distribuição das notas é positivo ( 0,41 ), logo a distribuição possui assimetria à direita.
As distribuições por Grande Região também apresentam assimetria positiva, concentração pouco maior do lado esquerdo do histograma e mais espalhada do lado direito.
Os Gráficos 3.2, 3.3 e 3.4 apresentam informações referentes à média da nota final dos Participantes, desagregando os resultados de acordo com, respectivamente, as Grandes Regiões do país, a Categoria Administrativa e a Organização Acadêmica.
Os gráficos apresentam o valor da média das notas como uma barra e os extremos do intervalo de confiança de 95 \% como linhas verticais unidas por uma linha horizontal na forma da letra H maiúscula.
Considerando -se o gráfico de notas segundo Grande Região ( Gráfico 3.2 ), observa -se que existe diferença estatisticamente significativa ao nível de 95 \% apenas entre a média obtida na região Sudeste ( 32,4 ) e as médias obtidas nas regiões Nordeste ( 34,3 ) e Sul ( 33,2 ).
Levando -se em conta os estudantes por Categorias Administrativas ( Gráfico 3.3 ), observa -se que existe diferença estatisticamente significante entre as médias das notas das IES Públicas e Privadas.
Pode ser observado ainda que a média dos alunos de IES Públicas ( 38,5 ) é maior que a dos alunos de IES Privadas ( 32,7 ).
A diferença entre IES Públicas e Privadas ( 5,8 ) é maior do que as diferenças entre regiões ( 3,4 ).
Tendo como foco o Gráfico 3.4, que apresenta as notas médias das provas segundo Organização Acadêmica, constata -se que existe diferença estatisticamente significativa ao nível de 95 \% nas médias das notas dos estudantes provenientes de Universidades, em relação aos provenientes de Centros Universitários e Faculdades.
A maior média foi obtida pelos estudantes de Universidades ( 33,6 ), e a menor, pelos de Centros Universitários ( 32,2 ). 3.1.2 Estatísticas Básicas no Componente de Formação Geral A Tabela 3.2 apresenta as estatísticas básicas em relação ao componente da prova que avalia a Formação Geral dos estudantes concluintes.
Os alunos de todo Brasil obtiveram desempenho médio de 40,7.
Quanto à variabilidade, o desvio padrão das notas dos estudantes do Brasil como um todo foi 15,7.
A maior média foi obtida na região Sul ( 41,9 ), e a menor, na região Norte ( 39,7 ).
As demais médias foram : 41,5 na região Centro-Oeste, 40,5 na região Nordeste e 40,2 na região Sudeste.
Já o maior desvio padrão foi obtido na região Centro-Oeste ( 17,0 ) e o menor na região Norte ( 14,8 ).
Os demais desvios padrões foram : 15,8 na região Sul, 15,7 na região Sudeste e 15,2 na região Nordeste.
A maior nota no componente de Formação Geral da prova do ENADE foi obtida por pelo menos um aluno nas regiões Sudeste e Sul ( 91,0 ) enquanto que a menor nota máxima foi obtida nas regiões Norte e Centro-Oeste ( 77,5 ).
Na região Nordeste a nota máxima foi 83,5.
A mediana do Brasil como um todo foi 41,0, sendo a menor mediana encontrada na região Nordeste ( 40,0 ) e a maior encontrada na região Sul ( 42,5 ).
A nota mínima nesta parte foi zero em quase todas as regiões, com exceção da região Norte ( 13,0 ).
O Gráfico 3.5 propicia a avaliação do desempenho dos estudantes no componente de Formação Geral a partir do histograma da distribuição das notas correspondentes.
A distribuição é unimodal, com moda em ( 40 ; 50 ], enquanto na prova como um todo a moda foi alcançada no intervalo ( 30 ; 40 ].
Nota -se, ainda, que no gráfico 3.5 as notas apresentam uma maior dispersão do que no Gráfico 3.1 ( distribuição das notas da prova ), confirmado pela comparação dos desvios padrões : 11,9 para a nota da prova como um todo e 15,7 para o componente de Formação Geral.
Para o componente de Formação Geral, o coeficiente de assimetria da distribuição das notas dos estudantes, ao contrário da prova como um todo, é negativo ( -0,09 ).
Por ser um coeficiente pequeno, a distribuição apresenta uma leve concentração à direita e cauda maior à esquerda.
Já nas regiões Norte e Nordeste os histogramas possuem assimetria levemente positiva ( respectivamente 0,27 e 0,03 ).
O Gráfico 3.5 propicia a avaliação do desempenho dos estudantes no componente de Formação Geral a partir do histograma da distribuição das notas correspondentes.
A distribuição é unimodal, com moda em ( 40 ; 50 ], enquanto na prova como um todo a moda foi alcançada no intervalo ( 30 ; 40 ].
Nota -se, ainda, que no gráfico 3.5 as notas apresentam uma maior dispersão do que no Gráfico 3.1 ( distribuição das notas da prova ), confirmado pela comparação dos desvios padrões : 11,9 para a nota da prova como um todo e 15,7 para o componente de Formação Geral.
Para o componente de Formação Geral, o coeficiente de assimetria da distribuição das notas dos estudantes, ao contrário da prova como um todo, é negativo ( -0,09 ).
Por ser um coeficiente pequeno, a distribuição apresenta uma leve concentração à direita e cauda maior à esquerda.
Já nas regiões Norte e Nordeste os histogramas possuem assimetria levemente positiva ( respectivamente 0,27 e 0,03 ).
Nos Gráficos 3.6, 3.7 e 3.8 são apresentadas as informações referentes ao desempenho dos Concluintes no componente de Formação Geral, em diferentes agregações : Grande Região do país, Categoria Administrativa e Organização Acadêmica.
Observa -se pelo Gráfico 3.6 que não existe diferença estatisticamente significativa entre a menor e a maior média das notas no Componente de Formação Geral, obtidas nas regiões Norte ( 39,7 ) e Sul ( 41,9 ).
Vemos que o intervalo de confiança mais largo é o da região Norte ; já o intervalo mais estreito é observado na região Sudeste.
Este fato está relacionado, também, com o tamanho da população envolvida, menor na região Norte do que na Sudeste.
No Gráfico 3.7, que representa as notas médias no Componente de Formação Geral segundo Categoria Administrativa do país, observa -se que existe diferença estatisticamente significativa entre as médias.
Assim como ocorreu para as notas da prova como um todo, para o componente de Formação Geral, os concluintes de Tecnologia em Gestão Financeira, das IES Públicas ( 46,8 ) obtiveram uma média maior do que os das IES Privadas ( 40,7 ).
Considerando -se o tipo de Organização Acadêmica, nota -se, no Gráfico 3.8, uma diferença estatisticamente significativa entre a maior e a menor média.
Nas Universidades ( 41,6 ), a média é mais elevada do que em Centros Universitários ( 41,0 ) e Faculdades ( 39,5 ). 3.1.3 Estatísticas Básicas do Componente de Conhecimento Específico A Tabela 3.3 apresenta as estatísticas básicas referentes ao componente de Conhecimento Específico da área de Tecnologia em Gestão Financeira.
A média do desempenho dos alunos do Brasil como um todo foi 30,1.
A maior média foi obtida na região Nordeste ( 32,2 ), e a menor, na região Norte ( 27,9 ).
As demais médias foram : 30,3 na região Sul e 29,8 nas regiões Sudeste e Centro-Oeste.
Quanto à variabilidade das notas, o desvio padrão do Brasil como um todo foi 13,3, sendo o maior desvio padrão observado na região Centro-Oeste ( 15,1 ) e o menor na região Norte ( 12,3 ).
Os demais desvios foram : 14,0 da região Nordeste, 13,8 da região Sul e 13,0 da região Sudeste.
A mediana das notas dos estudantes de todo o Brasil foi 28,6.
A maior mediana ocorreu na região Nordeste ( 30,5 ) e a menor na região Norte ( 25,0 ).
As demais medianas foram 28,6 na região Sul e 28,3 nas regiões Sudeste e Centro-Oeste.
A nota máxima do Brasil como um todo foi 91,0, sendo obtida por pelo menos um aluno da região Sul.
As demais notas máximas foram : 80,6 na região Sudeste, 77,8 na região Nordeste, 77,1 na Assim como os Gráficos 3.1 e 3.5, o Gráfico 3.9, apresentado a seguir, proporciona uma avaliação do desempenho de concluintes em relação ao componente de Conhecimento Específico com um histograma da distribuição das notas correspondentes.
Dentre as 3 distribuições apresentadas, esta é a mais concentrada nas notas baixas.
Esta também é uma distribuição unimodal, e o grupo modal é o ( 20 ; 30 ].
O coeficiente de assimetria da distribuição das notas no componente de Conhecimento Específico é positivo ( 0,55 ).
Nota -se pelo histograma ( Gráfico 3.9 ) que esta distribuição é concentrada nas notas mais baixas e possui cauda mais pesada à direita.
Em todas as regiões do país o coeficiente de assimetria também é positivo chegando a ser maior que 1,00 na região Norte ( 1,17 ). região Centro-Oeste e 73,6 na região Norte.
A nota mínima foi zero em quase todas as regiões, com exceção da região Norte, cuja nota mínima foi 9,4.
Assim como os Gráficos 3.1 e 3.5, o Gráfico 3.9, apresentado a seguir, proporciona uma avaliação do desempenho de concluintes em relação ao componente de Conhecimento Específico com um histograma da distribuição das notas correspondentes.
Dentre as 3 distribuições apresentadas, esta é a mais concentrada nas notas baixas.
Esta também é uma distribuição unimodal, e o grupo modal é o ( 20 ; 30 ].
O coeficiente de assimetria da distribuição das notas no componente de Conhecimento Específico é positivo ( 0,55 ).
Nota -se pelo histograma ( Gráfico 3.9 ) que esta distribuição é concentrada nas notas mais baixas e possui cauda mais pesada à direita.
Em todas as regiões do país o coeficiente de assimetria também é positivo chegando a ser maior que 1,00 na região Norte ( 1,17 ).
Os Gráficos 3.10, 3.11 e 3.12 apresentam, respectivamente, uma comparação dos resultados em relação à Grande Região do país, à Categoria Administrativa e à Organização Acadêmica, agora levando em conta o desempenho de alunos no componente de Conhecimento Específico da prova.
Pelo Gráfico 3.10, observa -se que existe diferença estatisticamente significativa entre as médias das notas no componente de Conhecimento Específico, da região Nordeste ( 32,2 ) em relação às regiões Sudeste ( 29,8 ) e Sul ( 30,3 ), devido estas regiões possuírem intervalo de confiança mais estreito.
Por outro lado, em relação às regiões Norte ( 27,9 ) e Centro-Oeste ( 29,8 ), devido estas regiões possuírem intervalo de confiança mais amplo, não existe diferença estatisticamente significativa.
A região Nordeste foi a que obteve média mais elevada, enquanto a região Norte foi a região que obteve média mais baixa ( 27,9 ).
Quanto à Categoria Administrativa ( Gráfico 3.11 ), observa -se um comportamento semelhante àquele da parte de Formação Geral e à prova como um todo, ou seja, existe diferença estatisticamente significativa entre as médias das IES Públicas e Privadas, sendo que a maior média foi obtida por alunos de IES Públicas de ensino ( 35,7 ).
Quanto ao Gráfico 3.12, observa -se, mais uma vez, que existe diferença estatisticamente significativa ao nível de 95 \% entre a maior e menor média das notas no componente de Conhecimento Específico.
Sendo que a média dos Concluintes das Universidades ( 30,9 ) foi maior do que de Centros Universitários ( 29,3 ) e de Faculdades ( 30,0 ). 3.2 ANÁLISE DAS QUESTÕES OBJETIVAS 3.2.1 Componente de Formação Geral A Tabela 3.4 apresenta as estatísticas básicas relativas às oito questões objetivas do componente da prova que abrange a Formação Geral dos estudantes.
A média do Brasil foi 44,7.
A menor média foi encontrada na região Norte ( 40,0 ) e a maior na região Sul ( 46,9 ).
As demais médias foram 46,0 na região Centro-Oeste, 44,9 na região Nordeste e 43,6 na região Sudeste.
O desvio padrão do Brasil foi 19,8, sendo o maior desvio padrão encontrado na região Centro-Oeste ( 21,6 ) e o menor na região Norte ( 17,9 ).
Os demais desvios foram : 20,0 na região Sudeste, 19,4 na região Sul e 18,9 na região Nordeste.
A mediana do Brasil como um todo ( 50,0 ) foi a mesma observada nas regiões Nordeste, Sul e Centro-Oeste, enquanto que nas demais regiões a mediana foi 37,5.
As notas máximas, nas questões objetivas do componente de Formação Geral, foram 100,0 em todas as regiões, com exceção da região Norte, cuja máxima foi 87,5.
Já as notas mínimas ( 0,0 ) foram iguais para todas as regiões.
A Tabela 3.5 apresenta os índices de facilidade e discriminação ( ponto bisserial ) das questões objetivas do componente de Formação Geral.
Quanto ao índice de facilidade, foram usadas as seguintes cores para diferenciar o nível de dificuldade da questão : • Azul para as questões classificadas com índice muito fácil ( índice \&gt ; =0,86 ), verde para as questões classificadas com índice fácil ( 0,61 a 0,85 ), amarelo para as questões classificadas com médio ( 0,41 a 0,60 ), vermelho para as questões classificadas com difícil ( 0,16 a 0,40 ) e roxo para as questões classificadas com muito difícil ( \&lt ; =0,15 ).
Já quanto ao índice de discriminação, foram usadas as seguintes cores para qualificar a questão : • As questões classificadas com índice fraco receberam a cor roxa ( índice \&lt ; =0,19 ), as classificadas com médio receberam a cor vermelho ( 0,20 a 0,29 ), as classificadas com bom receberam a cor verde ( 0,30 a 0,39 ) e as classificadas com muito bom ( \&gt ; =0,40 ) receberam a cor azul.
As questões objetivas do componente de Formação Geral, segundo o índice de facilidade, foram assim avaliadas : das oito questões, nenhuma teve o índice de facilidade classificado como fácil ou muito fácil.
Cinco questões foram consideradas de dificuldade média, situando -se no intervalo entre 0,41 e 0,60 do índice de facilidade, ou seja, houve entre 41,0 \% e 60,0 \% de acertos.
As demais questões foram classificadas como difíceis, situando -se no intervalo entre 0,16 e 0,40.
Como já comentado, para análise das questões objetivas relativas à Formação Geral segundo o poder de discriminação, utilizou -se, o índice de discriminação ponto bisserial.
Nesta análise as questões foram assim avaliadas : cinco das oito questões apresentaram 3.2 ANÁLISE DAS QUESTÕES OBJETIVAS 3.2.1 Componente de Formação Geral A Tabela 3.4 apresenta as estatísticas básicas relativas às oito questões objetivas do componente da prova que abrange a Formação Geral dos estudantes.
A média do Brasil foi 44,7.
A menor média foi encontrada na região Norte ( 40,0 ) e a maior na região Sul ( 46,9 ).
As demais médias foram 46,0 na região Centro-Oeste, 44,9 na região Nordeste e 43,6 na região Sudeste.
O desvio padrão do Brasil foi 19,8, sendo o maior desvio padrão encontrado na região Centro-Oeste ( 21,6 ) e o menor na região Norte ( 17,9 ).
Os demais desvios foram : 20,0 na região Sudeste, 19,4 na região Sul e 18,9 na região Nordeste.
A mediana do Brasil como um todo ( 50,0 ) foi a mesma observada nas regiões Nordeste, Sul e Centro-Oeste, enquanto que nas demais regiões a mediana foi 37,5.
As notas máximas, nas questões objetivas do componente de Formação Geral, foram 100,0 em todas as regiões, com exceção da região Norte, cuja máxima foi 87,5.
Já as notas mínimas ( 0,0 ) foram iguais para todas as regiões.
A Tabela 3.5 apresenta os índices de facilidade e discriminação ( ponto bisserial ) das questões objetivas do componente de Formação Geral.
Quanto ao índice de facilidade, foram usadas as seguintes cores para diferenciar o nível de dificuldade da questão : • Azul para as questões classificadas com índice muito fácil ( índice \&gt ; =0,86 ), verde para as questões classificadas com índice fácil ( 0,61 a 0,85 ), amarelo para as questões classificadas com médio ( 0,41 a 0,60 ), vermelho para as questões classificadas com difícil ( 0,16 a 0,40 ) e roxo para as questões classificadas com muito difícil ( \&lt ; =0,15 ).
Já quanto ao índice de discriminação, foram usadas as seguintes cores para qualificar a questão : • As questões classificadas com índice fraco receberam a cor roxa ( índice \&lt ; =0,19 ), as classificadas com médio receberam a cor vermelho ( 0,20 a 0,29 ), as classificadas com bom receberam a cor verde ( 0,30 a 0,39 ) e as classificadas com muito bom ( \&gt ; =0,40 ) receberam a cor azul.
As questões objetivas do componente de Formação Geral, segundo o índice de facilidade, foram assim avaliadas : das oito questões, nenhuma teve o índice de facilidade classificado como fácil ou muito fácil.
Cinco questões foram consideradas de dificuldade média, situando -se no intervalo entre 0,41 e 0,60 do índice de facilidade, ou seja, houve entre 41,0 \% e 60,0 \% de acertos.
As demais questões foram classificadas como difíceis, situando -se no intervalo entre 0,16 e 0,40.
Como já comentado, para análise das questões objetivas relativas à Formação Geral segundo o poder de discriminação, utilizou -se, o índice de discriminação ponto bisserial.
Nesta análise as questões foram assim avaliadas : cinco das oito questões apresentaram O Gráfico 3.13, para exemplificar, analisa o comportamento da questão de número 7 de Formação Geral.
Trata -se de uma questão de facilidade média e a que obteve o maior índice de discriminação dessa parte da prova.
Neste gráfico, cada uma das cinco curvas representa o percentual de respostas em determinada alternativa da questão, em função da nota dos estudantes nesta parte da prova ( Formação Geral/Múltipla Escolha ), antes de possíveis eliminações pelo critério do ponto bisserial.
A curva em verde corresponde à alternativa B, a correta para esta questão.
Assim, observa -se que entre os estudantes com notas mais baixas, nessa parte do exame, a situação mais frequente foi a escolha de uma das alternativas incorretas : a alternativa C ( em preto ) ou D ( em roxo ).
Na medida em que a nota aumenta, indicando desempenho melhor índices acima de 0,40 e, assim, foram classificadas com índice muito bom para esse grupo de alunos ; as outras três questões tiveram bom índice de discriminação, entre 0,30 e 0,39.
O índice de facilidade variou de 0,35 a 0,57, e o de discriminação, de 0,31 a 0,46.
As questões de número 1, 2, 6 e 7, classificadas com índice de discriminação muito bom, foram as mais fáceis desse conjunto : todas classificadas na categoria médio do índice de facilidade.
Em particular, a questão 6 foi a que apresentou o segundo maior poder discriminatório, com índice 0,43 e foi a segunda mais fácil.
A questão de número 5 apresentou índice de facilidade 0,35, ou seja, foi a que teve o menor quantitativo de acertos, dentre as questões objetivas relativas à Formação Geral.
Além disso, seu índice de discriminação foi bom, o que não a eliminou do cômputo da nota final.
A questão 7 obteve índice de discriminação muito bom, 0,46, e seu índice de facilidade foi médio.
Já a questão de número 8, considerada difícil, obteve índice de discriminação 0,36, considerado bom.
O Gráfico 3.13, para exemplificar, analisa o comportamento da questão de número 7 de Formação Geral.
Trata -se de uma questão de facilidade média e a que obteve o maior índice de discriminação dessa parte da prova.
Neste gráfico, cada uma das cinco curvas representa o percentual de respostas em determinada alternativa da questão, em função da nota dos estudantes nesta parte da prova ( Formação Geral/Múltipla Escolha ), antes de possíveis eliminações pelo critério do ponto bisserial.
A curva em verde corresponde à alternativa B, a correta para esta questão.
Assim, observa -se que entre os estudantes com notas mais baixas, nessa parte do exame, a situação mais frequente foi a escolha de uma das alternativas incorretas : a alternativa C ( em preto ) ou D ( em roxo ).
Na medida em que a nota aumenta, indicando desempenho melhor nesta parte da prova, aumenta concomitantemente a proporção de estudantes que selecionaram a alternativa correta B, atingindo 100 \% para os alunos com 8 acertos.
Essa análise permite verificar como a questão discriminou os grupos de desempenho, justificando o alto índice obtido na questão.
Os gráficos relativos às demais questões de Formação Geral constam do Anexo I. 3.2.2 Componente de Conhecimento Específico A Tabela 3.6 apresenta as estatísticas básicas em relação às questões objetivas do componente de Conhecimento Específico da prova, por Grande Região.
A média do Brasil deste componente foi de 34,2.
A menor média foi observada na região Norte ( 31,8 ) e a maior na região Nordeste ( 36,7 ).
O desvio padrão de todo o Brasil foi 15,1, sendo o menor desvio padrão encontrado na região Norte ( 14,2 ) e o maior na região Centro-Oeste ( 17,1 ).
A mediana de todo o Brasil foi 33,3, mesmo valor observado em todas as regiões do Brasil, exceto a região Norte ( 27,8 ).
A nota máxima da prova ( 100,0 ) foi obtida, nas questões objetivas do componente de Conhecimento Específico, por pelo menos um aluno da região A Tabela 3.7 apresenta os índices de facilidade e discriminação ( ponto bisserial ) das questões objetivas do componente de Conhecimento Específico.
Para facilitar a diferenciação das questões usaram -se as mesmas cores da Tabela 3.5 para as diferentes classificações dos índices de facilidade e de discriminação.
A partir dos índices obtidos das questões objetivas da parte da prova relativa ao componente de Conhecimento Específico, pode -se concluir que a maioria das questões objetivas da prova foi considerada pelo menos difícil : das 27 questões, dezenove foram classificadas como difíceis e três como muito difíceis.
Não houve questão classificada como fácil ou muito fácil, e outras cinco foram consideradas médias, entre 0,41 e 0,60.
Já quanto aos índices de discriminação das questões objetivas do componente de Conhecimento Específico da prova, tem -se como resultado a seguinte classificação : oito das 27 questões válidas foram consideradas como boas, enquanto não houve questão com índice de discriminação muito bom.
Dentre as demais, dez delas foram classificadas como médias e outras nove como fracas, sendo dezenove, por conseguinte, a quantidade de questões nos dois patamares mais baixos de discriminação.
Constata -se, assim, que a prova – no que se refere ao componente de Conhecimento Específico – possuía capacidade fraca de discriminar entre aqueles que dominam ou não o conteúdo.
Dentre as questões que alcançaram os maiores índices de discriminação, as de números 9, 11, 13, 14, 15, 17, 29 e 32 classificadas com índice bom, situando -se no intervalo de 0,30 a 0,39 do índice, quatro delas ( questões 11, 14, 17 e 29 ) foram Sul.
As notas máximas nas demais regiões foram : 88,9 nas regiões Nordeste, Sudeste e Centro-Oeste, e 83,3 na região Norte.
A nota mínima foi zero em quase todas as regiões do Brasil, sendo a região Norte ( 11,1 ), a exceção.
A Tabela 3.7 apresenta os índices de facilidade e discriminação ( ponto bisserial ) das questões objetivas do componente de Conhecimento Específico.
Para facilitar a diferenciação das questões usaram -se as mesmas cores da Tabela 3.5 para as diferentes classificações dos índices de facilidade e de discriminação.
A partir dos índices obtidos das questões objetivas da parte da prova relativa ao componente de Conhecimento Específico, pode -se concluir que a maioria das questões objetivas da prova foi considerada pelo menos difícil : das 27 questões, dezenove foram classificadas como difíceis e três como muito difíceis.
Não houve questão classificada como fácil ou muito fácil, e outras cinco foram consideradas médias, entre 0,41 e 0,60.
Já quanto aos índices de discriminação das questões objetivas do componente de Conhecimento Específico da prova, tem -se como resultado a seguinte classificação : oito das 27 questões válidas foram consideradas como boas, enquanto não houve questão com índice de discriminação muito bom.
Dentre as demais, dez delas foram classificadas como médias e outras nove como fracas, sendo dezenove, por conseguinte, a quantidade de questões nos dois patamares mais baixos de discriminação.
Constata -se, assim, que a prova – no que se refere ao componente de Conhecimento Específico – possuía capacidade fraca de discriminar entre aqueles que dominam ou não o conteúdo.
Dentre as questões que alcançaram os maiores índices de discriminação, as de números 9, 11, 13, 14, 15, 17, 29 e 32 classificadas com índice bom, situando -se no intervalo de 0,30 a 0,39 do índice, quatro delas ( questões 11, 14, 17 e 29 ) foram classificadas na categoria média, e as outras quatro como difíceis, quanto ao índice de facilidade.
A questão de número 20 foi a mais difícil dentre as 27 questões específicas, com baixo índice de facilidade, apenas 11,0 \% de acertos.
Essa questão apresentou poder discriminatório igualmente baixo, 0,12, o que comprova ter sido esta a mais difícil para os estudantes.
Destaca -se, também, a questão 23, com índice de facilidade 0,13, o que, em termos percentuais, corresponde a 13,0 \% de estudantes que responderam acertadamente, obtendo, ainda, 0,07 de índice de discriminação.
Tais questões foram, portanto, duas das mais difíceis da prova.
Pelo critério do ponto bisserial as questões 20 e 23 foram eliminadas do cômputo da nota final.
Além destas duas, as demais questões com índice fraco de discriminação, questões 12, 16, 19, 21, 26, 27 e 30 também não foram computadas.
A título de exemplo das análises do comportamento das questões objetivas, o Gráfico 3.14 analisa a questão 11 do componente de Conhecimento Específico.
Esta foi a questão mais fácil da prova, apresentando índice de facilidade 0,57, ou seja, 57,0 \% dos estudantes assinalaram acertadamente a opção D, correspondente ao gabarito.
Seu índice de discriminação foi igual a 0,38, classificado como bom.
Neste gráfico, cada uma das cinco curvas representa o percentual de respostas em determinada alternativa da questão 11, em função da nota dos estudantes nesta parte da prova, antes de possíveis eliminações de questões pelo critério do ponto bisserial.
A alternativa correta D, representada no gráfico pela curva em roxo, foi escolhida em maiores proporções pelos alunos com desempenho melhor nesta parte da prova.
Já as alternativas incorretas, também denominadas distratores, foram selecionadas principalmente por Aqueles com nota zero, na sua quase totalidade deixaram esta questão em branco ou marcaram mais de uma alternativa, comportamento considerado inválido.
A proporção de alunos que selecionou a resposta correta D aumenta gradativamente, chegando a atingir 100 \% para alunos com 18 acertos ou mais, enquanto a proporção dos que escolheram alternativas incorretas decai, a partir da primeira nota não nula, como função da nota nesta parte da prova.
Os gráficos relativos às demais questões do componente de Conhecimento Específico constam do Anexo I. 3.3 ANÁLISE DAS QUESTÕES DISCURSIVAS 3.3.1 Componente de Formação Geral As análises dos resultados de desempenho dos estudantes de Tecnologia em Gestão Financeira nas duas questões discursivas relativas à Formação Geral encontram -se na Tabela 3.8 e no Gráfico 3.15.
Na tabela 3.8 observa -se que as notas médias foram mais baixas nesse conjunto de questões do que no das objetivas.
Os estudantes de todo o Brasil obtiveram, em Formação Geral, média 44,7 nas questões objetivas e 34,9 nas questões discursivas.
A mediana confirma o pior desempenho dos alunos de todo o Brasil nas questões discursivas do componente de Formação Geral.
Enquanto essa estatística foi de 37,5 para questões discursivas, para as questões objetivas essa estatística foi de 50,0.
Pode -se, também, notar um aumento do desvio padrão de 19,8, nas questões objetivas do componente de Formação Geral dos alunos de todo o Brasil, para 21,2 nas questões discursivas do mesmo componente.
Como já mencionado, a mediana de todo o Brasil, neste componente, foi 37,5, sendo o mesmo valor encontrado nas regiões Nordeste, Sudeste e Sul, enquanto a maior mediana foi encontrada na região Norte ( 43,8 ) e no Centro-Oeste essa estatística foi de 40,0.
A maior nota máxima ( 87,5 ) foi obtida, nas questões discursivas no componente de Formação Geral, por pelo menos um aluno da região Sul, enquanto a menor nota máxima ( 75,0 ) foi obtida na região Norte.
A nota mínima foi zero em todas as regiões do Brasil, sem exceção.
O Gráfico 3.15 representa a distribuição das notas nas questões discursivas no componente de Formação Geral.
A distribuição é bimodal com modas nos intervalos [ 0 ; 10 ], devido à grande quantidade de notas zero e à alta frequência de alunos que deixaram este tipo de questão em branco, e ( 40 ; 50 ].
A distribuição possui assimetria à esquerda, com coeficiente de assimetria -0,30.
Em todas as regiões o coeficiente de assimetria também é negativo, sendo o maior em módulo na região Norte ( -0,51 ), e o menor na região Nordeste ( -0,20 ).
Na sequência, os resultados verificados para cada uma das questões discursivas de Formação Geral serão apresentados, estabelecendo -se relações com os conteúdos abordados em cada uma delas.
Os comentários da Banca de docentes corretores a respeito do observado na correção das respostas dos estudantes, suas impressões e conclusões serão apresentados junto à análise de cada questão.
Cumpre esclarecer que, tendo em vista que as questões discursivas de Formação Geral são padronizadas, ou seja, constam de todas as provas, os comentários da Banca são os mesmos para todas as carreiras acadêmicas, sendo direcionados a todos os estudantes que participaram do ENADE/2012.
A seguir, serão analisados os desempenhos da Área de Tecnologia em Gestão Financeira nas duas questões discursivas de Formação Geral do ENADE/2012, comparando os resultados obtidos com comentários para cada questão. 3.3.1.1 Análise da Questão Discursiva 1 do Componente de Formação Geral Os dados de Tecnologia em Gestão Financeira, obtidos a partir das respostas à questão 1, encontram -se na Tabela 3.9 e no Gráfico 3.16.
Nessa questão – de pior desempenho dentre as duas de Formação Geral – os alunos de todo Brasil tiveram média, 29,1.
A maior média para a questão 1 foi obtida na região Norte ( 34,4 ), e a menor, na região Nordeste ( 27,4 ).
Quanto à variabilidade das notas, o desvio padrão de todo o Brasil foi 24,0.
O menor desvio padrão foi obtido na região Sudeste ( 23,7 ) e o maior desvio padrão foi obtido na região Centro-Oeste ( 25,7 ).
As medianas foram as mesmas para todas as regiões do Brasil ( 30,0 ), com exceção da região Norte cuja mediana foi 37,5.
A maior nota máxima da questão discursiva 1, 95,0, foi obtida nas regiões Sudeste e Sul, enquanto a menor nota máxima ( 80,0 ) foi obtida na região Norte.
Além disso, a mesma nota mínima ( 0,0 ) foi obtida em todas as regiões do país, sem exceção.
O Gráfico 3.16 mostra a distribuição das notas na questão discursiva 1 do componente de Formação Geral.
Observa -se que a maior frequência, correspondente à moda, ocorre na classe dos alunos que deixaram a questão em branco.
Considerando apenas os estudantes que tentaram resolver a questão, a moda ocorre no intervalo ( 20 ; 30 ].
Também observa -se grande concentração de notas no intervalo ( 30 ; 40 ].
A distribuição é assimétrica à direita, com coeficiente de assimetria positivo ( 0,29 ).
Em todas as regiões do Brasil este coeficiente se mantém positivo ( entre 0,04 na região Norte e 0,41 na região Nordeste ). 3.3.1.2 Comentários sobre a correção das respostas à Questão Discursiva 1 A questão aborda um tema bastante polêmico o qual vem sendo amplamente discutido, a conciliação entre o desenvolvimento econômico e a preservação dos recursos naturais.
Busca -se observar a relação entre o estímulo governamental à compra de veículos novos e os impactos ambientais decorrentes desta política pública.
Sua compreensão implica no conhecimento de uma temática atual e na reflexão critica a respeito das decisões governamentais em sua correlação com o cenário econômico internacional e os problemas e soluções de ordem ambiental.
Embora o comando da questão solicitasse explicitamente que a resposta fosse fornecida em um texto único ( “ redija um texto dissertativo sobre... ” ), foi significativamente alta a proporção de respondentes que atendeu a cada item ( a, b e c ) do enunciado de maneira isolada, como partes estanques da mesma resposta.
No que se segue encontra -se uma avaliação da correção de cada um dos três itens. a - “ conceito de desenvolvimento sustentável ” A variedade de respostas a esse item confirmou a importância do tema para avaliar a formação geral dos estudantes e cursos, já que ( i ) é atual ; ( ii ) foi abordado nos meios de comunicação ( especialmente no ano de 2012, em que o Brasil sediou a “ Rio+20 ” ) ; ( iii ) é relevante não apenas no Brasil como em todos os principais fóruns internacionais ; e ( iv ) é comum a diversas áreas de conhecimento.
Outros dois aspectos positivos merecem destaque : ( i ) a grande incidência de respostas que apontaram a relação entre gerações como aspecto essencial da ideia de sustentabilidade socioambiental – no mais, as respostas apresentaram pelo menos a tentativa de indicar uma definição de desenvolvimento sustentável ou a abordagem positiva dessa mesma noção, evidenciado compreensão do enunciado ; e ( ii ) o exercício crítico que a questão propunha, no sentido de exigir dos respondentes a compreensão das contradições de políticas públicas e a apresentação de propostas de soluções para conciliar o incentivo do crescimento econômico e a promoção da sustentabilidade socioambiental, a partir de uma perspectiva transformadora da realidade.
Apesar dos aspectos positivos apontados no parágrafo anterior, a maioria dos estudantes não alcançou o máximo da pontuação.
De forma geral, trataram o desenvolvimento sustentável como sinônimo de não agressão ambiental.
Alguns confundiram desenvolvimento sustentável com desenvolvimento econômico, e outros faziam referência apenas à preocupação com a conservação dos recursos mencionando a preocupação com gerações futuras.
Deve -se, ainda, destacar o significativo número de alunos que apresentou uma definição aparentemente copiada do enunciado de outra questão da prova, como por exemplo, a questão objetiva 31 da prova de Administração. b - “ conflito entre o estímulo à compra de veículos automotores e a promoção da sustentabilidade ” Esse foi o item identificado como aquele no qual as respostas dos estudantes avaliados mais se aproximaram da pontuação máxima prevista no padrão de resposta.
Ou seja, a ampla maioria dos respondentes foi bem-sucedida em apontar a relação entre o aumento da quantidade de veículos automotores em circulação e os danos socioambientais que isso acarreta no cenário urbano, devido ao aumento da emissão de gases poluentes.
No item b, as respostas corretas se concentraram em críticas ao incentivo oferecido pelo governo, referências ao problema dos congestionamentos e ao aumento da poluição.
Contudo, a maioria das respostas não explicitava a correlação entre o aumento da circulação de veículos automotores e o agravamento dos problemas ambientais.
Um número expressivo de alunos apresentou uma resposta em que era dado destaque à política de redução do IPI como um mecanismo para renovação da frota de automóveis ou para a aquisição do transporte próprio e da consequente redução da dependência ao transporte coletivo.
Também foram notados alguns tipos de padrões recorrentes nas tentativas de resposta a este item.
Dois desses padrões se destacaram : ( i ) as declarações genéricas de condenação do “ consumismo ” de uma “ nova classe média ” ou da busca pelo “ status ” que a posse de um automóvel particular proporcionaria, como se fosse possível atribuir a essa ordem de coisas a causa principal do aumento dos congestionamentos nas grandes cidades ; e ( ii ) a repetição da avaliação, presente no enunciado da questão, que aponta de forma superficial a contradição entre o estímulo à demanda por automóveis particulares e a promoção da sustentabilidade. c - “ ações de fomento ao transporte urbano sustentável no Brasil ” Este foi o item que ofereceu maior facilidade de aplicação do padrão de resposta, mas foram poucos os casos de pontuação integral.
Dentre as soluções mais destacadas pelos alunos apareceram : ( i ) o incentivo ao uso de energias alternativas ; e ( ii ) o estímulo ao transporte coletivo ( metrô, trem e ônibus ), além do uso de bicicletas.
Importante destacar o fato de que muitos apresentavam respostas com referências à necessidade de incentivo ao transporte urbano, nas quais este termo era utilizado como sinônimo de transporte coletivo.
A avaliação da questão 1 trouxe à tona problemas estruturais da educação básica, que se refletem no ensino superior.
Foi muito grande o número de provas em que o aluno demonstrou não compreender o enunciado da questão, apresentando respostas que fugiam ao tema proposto, mesmo quando bem redigidas.
Uma resposta recorrente fazia referência às consequências da ampla oferta de crédito à população e aos problemas relacionados à inadimplência, situação que, embora, verdadeira, não estava relacionada à questão proposta.
Houve, ainda, os casos em que as respostas apresentadas não só se afastaram do tema como não reuniam qualquer tipo de argumentação coerente.
E, por fim, destaca -se a existência de respostas incoerentes e desconexas.
Outro aspecto geral diz respeito ao grande número de provas em branco ou em que o aluno apresentou a resposta a outro item no espaço da questão 1.
Considerando -se o objetivo geral deste exame, avaliar a qualidade dos cursos de graduação oferecidos no país e contribuir para a melhoria da qualidade do ensino superior, destaca -se a imensa responsabilidade do processo que envolve a elaboração deste diagnóstico.
Nesse sentido, é urgente alertar para a ocorrência de um número significativo de respostas que se afastam de um padrão mediano de qualidade esperado de um aluno concluinte em um curso de graduação, especialmente em se considerando o fato de que a questão proposta faz referência a tema intensamente explorado pela mídia. 3.3.1.3 Análise da Questão Discursiva 2 do Componente de Formação Geral A Tabela 3.10 mostra que o desempenho dos estudantes na questão 2 ( média 40,6 ) foi tipicamente superior ao obtido na questão de número 1 ( média 29,1 ).
A região Norte foi aquela onde a média, nessa questão, foi maior ( 44,2 ), e a de menor média foi a região Centro-Oeste ( 38,9 ).
Quanto à variabilidade das notas, o desvio padrão de todo o Brasil foi 24,9, superior ao obtido na questão de número 1 ( 24,0 ).
O maior desvio nessa questão foi obtido na região Centro-Oeste ( 25,9 ), enquanto o menor foi obtido na região Sudeste ( 24,6 ).
Assim como no Brasil como um todo, a mediana foi de 45,0 nas regiões Sudeste, Sul e Centro-Oeste, enquanto essa estatística foi de 50,0 nas demais regiões.
A maior nota máxima ( 95,0 ) foi obtida nas Sudeste e Sul, enquanto a menor nota máxima ( 85,0 ) foi obtida nas regiões Norte e Centro-Oeste.
A nota mínima foi zero em todas as regiões do Brasil, sem exceção.
O Gráfico 3.17 mostra a distribuição das notas na questão discursiva 2 do componente de Formação Geral.
Neste gráfico, apesar da grande quantidade de alunos que deixaram a questão 2 em branco ( em torno de 19 \% ), observa -se a moda principal no intervalo ( 40 ; 50 ].
O coeficiente de assimetria das notas nesta questão é negativo ( -0,47 ), indicando que a distribuição possui assimetria à esquerda.
O mesmo ocorre para todas as regiões do país, cujos coeficientes variam de -0,64 na região Norte a -0,40 na região Sul. 3.3.1.4 Comentários sobre a correção das respostas à Questão Discursiva 2 O conteúdo da questão era bastante adequado às Diretrizes do INEP para o ENADE, além de atual e de abrangente.
Porém, exigia -se a articulação de muitas variáveis.
Buscando aferir, a partir da apresentação de uma charge e de um texto introdutório, as competências e habilidades relacionadas à leitura, à interpretação, à análise crítica, à capacidade argumentativa, à elaboração de sínteses e de avaliação prospectiva sobre tema da atualidade, a questão tratava da violência e suas repercussões no contexto atual.
O comando da questão solicitava, a partir da definição de violência da OMS e de duas imagens, que o aluno elaborasse um texto, de caráter dissertativo, sobre a “ violência na atualidade ”.
Embora tal solicitação fosse explícita ( “ redija um texto dissertativo sobre... ” ), como na questão 1, foi alta a proporção de alunos que respondeu a cada item do enunciado isoladamente, como partes independentes.
Verificaram -se poucas respostas em branco, como também foram poucas as respostas contendo protestos contra o ENADE, muito frequentes em anos anteriores, o que significa um avanço para a afirmação do ENADE como um exame aceito socialmente. a -"tecnologia e violência"Tanto as charges quanto o texto apresentado no enunciado induziram os alunos a focarem suas respostas nos meios de comunicação.
Além disso, algumas respostas estabeleciam a tecnologia, em si, como causa da violência tanto nas escolas ( em particular ), quanto na sociedade ( no geral ).
Muitos respondentes tiveram dificuldade para fazer a relação entre tecnologia, violência e a escola, de modo adequado, centrando -se na violência nas escolas.
Além disso, a grande maioria das respostas girava em torno da exposição da violência através da TV, internet e jogos digitais, sendo que predominaram as abordagens com ênfase em programas e filmes veiculados pela televisão.
Poucos foram os estudantes que estabeleceram relação entre o contexto socioeconômico e a violência. b -"causas e consequências da violência na escola"Embora a solicitação desse item fosse objetiva e atendida por grande parte dos alunos, inúmeras vezes, observou -se confusão entre o significado de causa e de consequência.
Em outras inúmeras vezes, as respostas apontavam apenas para uma das classes, causas ou consequências.
Foram frequentes as respostas que atribuíam o aumento da violência à ineficiência do poder público, apontando a falta de investimentos, o baixo salário dos professores, a ausência de profissionais qualificados como causas.
Foram frequentes também os que apontavam a violência doméstica, o bullying, e a transferência do imaginário dos filmes e jogos para a realidade, como causas da violência. c -"proposta de solução para o problema da violência na escola"Sem dúvida, há uma grande diversidade de abordagens possíveis para apresentação de uma proposta de solução para o problema da violência, o que o padrão de resposta buscou contemplar.
No entanto, houve grande número de respostas baseadas em opiniões pessoais dos estudantes e vinculadas ao senso comum.
As respostas giraram em torno da necessidade de mais investimentos na educação por parte do poder público e seu rebatimento no maior aparelhamento das escolas, salário dos professores etc.
Além disso, foram muitas as respostas que trataram da necessidade de maior conscientização dos pais, alunos e professores com a realização de um trabalho conjunto que envolvesse todos os segmentos da sociedade.
De um modo geral, os estudantes demonstraram dificuldades em relação ao nível de abrangência exigido, falta de conteúdo, inabilidade na elaboração de texto escrito.
Da grande quantidade de provas corrigidas, o que ficou evidente, foi a dificuldade de estabelecer as conexões requeridas na questão 2.
A violência nas escolas foi bastante associada ao reflexo do meio onde a criança e o adolescente mantém os seus vínculos afetivos.
Surgiram muitas respostas revelando críticas à intolerância no núcleo familiar, à jornada de trabalho dos pais, ao distanciamento da mãe nas atividades domésticas etc.
Em boa parte das respostas, os estudantes lançaram a causa da violência escolar exclusivamente na ausência de participação familiar na vida dos filhos, às vezes criando uma cisão com a proposta da questão, já que abandonavam totalmente a temática da tecnologia em suas reflexões.
Além da ausência dos pais no cotidiano da escola, divórcios e atos de violência contra crianças foram considerados fatores que delegariam para a escola novas funções.
Também foram mencionadas a má formação de professores e a insatisfação com o ambiente escolar.
Em inúmeros casos, ficou evidente o anseio de se fazer da escola um lugar de resgate tanto da criança e do adolescente quanto da família, compreendendo -se a escola em sentido social e assistencial.
Segundo este amplo grupo de concluintes, a escola deveria oferecer -se como centro de cultura e lazer, como espaço para o acompanhamento psicológico e serviço social às comunidades diretamente envolvidas.
Ressalte -se o constante aparecimento das sugestões de horário integral, melhores salários para professores e iniciativas para melhoria na formação de professores.
As propostas de solução para acabar com a violência na escola se detiveram em aspectos muito gerais e poucos alunos detalharam uma proposta de intervenção, apresentando ações concretas.
Dentre as propostas de solução surgiram : educação familiar, ensino religioso, maior segurança nas escolas, com utilização de aparelhos para detecção de armas, e o retorno do papel do professor como um profissional respeitado pelos poderes municipais e estaduais.
Um percentual significativo de provas apresentou resposta simplista, com ênfase em determinados pontos como : religião ( sem demonstração do conceito de laicidade ), políticas sociais ( sem demonstração do domínio de conceitos que envolvem Estado, Governo e Politicas Públicas ) etc.
Foi grande a incidência de respostas que se referiam às leis, à vigilância e maior controle dos pais e professores ; no entanto, não demonstram entender a noção de legalidade para os diferentes segmentos sociais e como os mesmos se encontram submetidos a diferentes formas de poder.
Alguns alunos, ainda que raros, trataram a questão com a profundidade esperada de concluintes de cursos universitários.
Por outro lado, foram vários os textos que apresentaram uma argumentação simplória e desconexa, com forte visão paternalista sobre o papel do governo e do professor, como se tais instâncias pudessem intervir na educação familiar.
Enquanto alguns conseguiram desenvolver o tema proposto com elegância e desenvoltura, encadeando o tema de forma harmônica, com excelente nível de linguagem escrita e correção gramatical, outros sequer conseguiram conectar as ideias e construir um texto razoável, utilizando, inclusive, um linguajar mais propício à linguagem falada do que à escrita.
Cabe também ressaltar a quantidade de respostas carregadas de conteúdo com juízo de valor discriminatório ou extremamente conservador, do tipo : “ mulher não deve trabalhar fora ” ; “ amar e temer a Deus é a solução ” ; obrigatoriedade do ensino religioso e de uma doutrina religiosa ; imposição de uma educação do tipo militar nas escolas ; censura nos meios de comunicação etc.
Outras soluções apontadas para resolver o problema da violência nas escolas estavam associadas à, por exemplo : redução da maioridade penal, defesa da pena de morte, milícias armadas na porta das escolas, punições como o retorno da palmatória.
O Estatuto da Criança e do Adolescente foi uma peça legal criticada, e recebeu, em diversas respostas, o título de responsável pelo sentimento de impunidade de adolescentes infratores.
Contudo, algumas dessas proposições, desde que bem contextualizadas e formuladas com boa argumentação, foram pontuadas de forma a respeitar o direito democrático de pluralidade opinativa.
É imperioso ressaltar que a falta de inteligibilidade segue representando um aspecto merecedor de providências na formação universitária, uma vez que exprime uma habilidade importante para qualquer profissional de nível superior.
Além de terem sido poucas respostas boas e claras, foi possível identificar a deficiência de argumentação, além dos inúmeros problemas gramaticais ( ortografia, concordância, pontuação, regência, etc. ).
Salienta -se, nesse sentido, o uso indiscriminado de letras maiúsculas e minúsculas numa mesma frase, sem obedecer ao critério formal da língua portuguesa. 3.3.2 Componente de Conhecimento Específico Na parte da prova relativa às questões discursivas no componente de Conhecimento Específico ( Tabela 3.11 ), observa -se que a média foi bem mais baixa do que para as questões discursivas do componente de Formação Geral.
Enquanto no componente de 63 Formação Geral a média para estudantes de Tecnologia em Gestão Financeira de todo o Brasil foi 34,9, na parte de Conhecimento Específico a média foi 6,5.
A maior média deste componente foi obtida pelos estudantes da região Sudeste ( 6,6 ), e a menor, pelos das regiões Norte ( 6,0 ).
Quanto à variabilidade das notas, o desvio padrão de todo o Brasil foi 8,6.
O maior desvio padrão foi encontrado na região Centro-Oeste ( 9,1 ), e o menor, na região Norte ( 6,9 ).
A maior nota máxima, nas questões discursivas do componente de Conhecimento Específico, foi obtida na região Sudeste ( 76,7 ) enquanto a menor nota máxima foi encontrada na região Norte ( 30,0 ).
Além disso, a nota mínima ( 0,0 ) foi obtida por alunos de todas as regiões do Brasil sem exceção.
A mediana do Brasil como um todo foi 3,3, também encontrada nas regiões Sudeste e Sul.
As demais medianas foram : 5,0 na região Norte e 1,7 nas regiões Nordeste e Centro-Oeste.
O Gráfico 3.18 representa a distribuição das notas nas questões discursivas no componente de Conhecimento Específico.
Essa distribuição tem moda no intervalo de [ 0 ; 10 ].
Nota -se que a partir desse intervalo a frequência de notas nos intervalos é cada vez menor.
Assim, o coeficiente de assimetria é positivo ( 1,83 ).
Este coeficiente também é positivo e maior do que 1,00 em todas as regiões do país ( entre 1,26 na região Norte e 1,91 na região Sudeste ). 3.3.2.1 Análise da Questão Discursiva 3 do Componente de Conhecimento Específico Na questão 3, cujos resultados aferidos encontram -se descritos na Tabela 3.12, a média dos estudantes de todo o Brasil foi 7,0.
A menor média nessa questão foi obtida pelos alunos da região Nordeste e Centro-Oeste ( 6,7 ), enquanto a maior média foi obtida na região Norte ( 9,2 ).
Quanto à variabilidade das notas, o desvio padrão de todo o Brasil foi 13,6.
O maior desvio padrão foi obtido na região Sudeste ( 13,9 ), enquanto o menor foi obtido na região Sul ( 13,1 ).
A nota máxima, 100,0 pontos, foi alcançada por pelo menos um aluno da região Sudeste.
A região onde ocorreu a menor nota máxima foi a Norte ( 55,0 ).
A mediana e a nota mínima do Brasil como um todo foram zero, e foram as mesmas em todas as regiões do país, ou seja, pelo menos metade dos alunos tirou nota zero nesta questão.
O Gráfico 3.19 mostra a distribuição das notas na questão discursiva 3, do componente de Conhecimento Específico.
Essa distribuição é bimodal, com modas nas questões em branco e no intervalo ( 0 ; 10 ].
Observamos que o gráfico possui assimetria à direta, comprovada pelo coeficiente de assimetria 2,54.
Em todas as regiões do país também há coeficiente de assimetria positivo e alto ( entre 1,72 na região Norte e 2,94 na região Centro-Oeste ). 3.3.2.2 Comentários sobre as respostas à Questão Discursiva 3 A questão buscava avaliar o conhecimento do conceito de responsabilidade socioambiental e como a parte financeira da empresa era afetada pela adoção deste tipo de atuação.
Nessa questão, além da média ser baixa, foi bastante significativo que o número de notas 0,0 ( dentre os alunos que tentaram responder ) somado à quantidade dos que deixaram a questão em branco, respostas desconsideradas ou anuladas.
Este total correspondeu a mais do que 60 \% dos concluintes.
Tais alunos podem não terem sido expostos ao tema em seus cursos ou, caso contrário, não conseguiram aprender os conceitos exigidos na questão. 3.3.2.3 Análise da Questão Discursiva 4 do Componente de Conhecimento Específico A Tabela 3.13 contém as informações relativas à questão 4 do conjunto de questões do componente de Conhecimento Específico.
O desempenho dos estudantes de todo o Brasil nesta questão foi tipicamente superior ao desempenho na questão de número 3.
A média geral do Brasil foi 11,6, sendo a menor média registrada na região Norte ( 8,8 ) e a maior na região Sudeste ( 11,7 ).
A maior nota máxima ( 95,0 ) foi atingida nas regiões Sudeste e Sul, enquanto a menor nota máxima ( 55,0 ) foi obtida na região Norte.
A mediana foi 5,0 para o Brasil como um todo e também para as regiões Sudeste e Sul, enquanto nas regiões Norte, Nordeste e Centro-Oeste a mediana foi 0,0, ou seja, pelo menos metade dos estudantes dessas três regiões tirou nota zero nesta questão.
A nota mínima ( 0,0 ) foi obtida em todas as regiões do Brasil, sem exceção.
O Gráfico 3.20 representa a distribuição de notas na questão discursiva 4, no componente de Conhecimento Específico.
Como nas outras questões discursivas, o número de estudantes que deixou a questão em branco foi muito grande, correspondendo à moda da distribuição.
Dentre as notas dos que resolveram a questão 4, o intervalo modal foi ( 10 ; 20 ].
A assimetria da distribuição confirma o valor baixo da média, que foi 11,6, já que apresenta concentração em notas inferiores ao intervalo modal e uma cauda longa do lado direito deste intervalo.
Tanto na distribuição como um todo ( 1,52 ) quanto na análise dos resultados dos alunos por região ( entre 1,48 na região Sudeste e 1,65 na região Norte ), os coeficientes de assimetria são sempre positivos e com valor alto, sempre maiores do que 1,00. 3.3.2.4 Comentários sobre as respostas à Questão Discursiva 4 A questão 4 avaliava se o concluinte conhecia o conceito de spread bancário e sua relação a fatores macroeconômicos e financeiros.
Consideramos que essa questão abordava assunto mais intimamente ligado à Área de concentração do curso, se comparada à questão 3.
Provavelmente por isso, houve uma melhora significativa dos resultados, sugerindo que os temas específicos são mais tratados ou mais aprofundados nos cursos.
Apesar disso, foi também bastante significativo o percentual de notas 0,0 ( aluno que tentou responder ) somado ao de brancos, desconsiderados e nulos.
Assim, mais de 48 \% do total de notas atribuídas resultaram iguais a zero. 3.3.2.5 Análise da Questão Discursiva 5 do Componente de Conhecimento Específico A Tabela 3.14 contém as informações relativas à questão 5 do conjunto do componente de Conhecimento Específico.
O desempenho dos estudantes nessa questão foi inferior ao das questões 3 e 4.
A nota média dos estudantes de todo o Brasil foi 0,9.
A maior média foi registrada nas regiões Nordeste e Centro-Oeste ( 1,2 ), enquanto a menor média foi registrada na região Norte ( 0,0 ).
Quanto à variabilidade das notas, o desvio padrão dos alunos do Brasil, como um todo, foi 6,4.
Enquanto o maior desvio foi encontrado na região Nordeste ( 7,3 ), o menor foi encontrado na região Norte ( 0,0 ).
A nota máxima ( 60,0 ) foi alcançada em todas as regiões, com exceção da região Norte, sendo possível constatar que todos os alunos dessa região tiveram nota zero na questão.
A mediana e a nota mínima foram zero para todas as regiões do Brasil, ou seja, assim como na questão 4 do componente de Conhecimento Específico, pelo menos metade dos alunos tirou nota zero na questão.
O Gráfico 3.21 representa a distribuição das notas da questão discursiva 5 do componente de Conhecimento Específico.
Destaca -se, como nas demais questões discursivas, o grande número de estudantes que deixaram esta questão em branco.
Entre entre as três questões discursivas do Componente de Conhecimento Específico, esta foi a que apresentou um menor espalhamento das notas daqueles que resolveram a questão, também sendo relevante não terem sido registradas notas acima de 60,0.
O coeficiente de assimetria para todos os alunos é 7,71. 3.3.2.6 Comentários sobre as respostas à Questão Discursiva 5 A temática da questão 5 era bastante específica e de grande utilização prática para os profissionais de gestão financeira.
Todas as IES devem obrigatoriamente abordar o tema da questão por se tratar de item fundamental para o aluno.
Porém, devido a falhas no enunciado o padrão de resposta considerou mais do que uma alternativa de condução da resposta em relação à margem de lucro ( operacional ou líquida ).
O fato da segunda parte da questão depender do resultado correto da primeira pode ter prejudicado a pontuação final da questão.
De fato, mesmo sabendo responder à segunda parte, caso se o aluno errasse a primeira parte obrigatoriamente erraria a segunda.
A questão 5 foi a de maior quantidade de notas zero.
Se somadas a quantidade de zeros de alunos que responderam erroneamente com a dos que deixaram a questão em branco, tiveram a resposta desconsiderada ou anulada, obtém -se mais do que 92 \% do total de notas atribuídas.
Tal resultado pode estar influenciado pelo fato de ser uma questão discursiva que dependia do uso de dados quantitativos, além de ser a última da prova. 3.3.3 Considerações Finais O desempenho dos concluintes foi pior do que o esperado.
Causou impacto observar que um conceito básico na Área de finanças, como spread bancário, não ser dominado por cerca de 50 \% dos concluintes.
Preocupante também foi o baixo desempenho na questão 5, que exigia o uso de dados quantitativos.
Provavelmente, este é um resultado do desempenho muito baixo em Matemática na Educação Básica, evidenciado por diversas avaliações de grande porte.

% matched lemmas: 
\end{textsample}
